% =====================================================================
%  LTL CLASSROOM OBSERVATION MANUAL
%  Laid out in the format of the original TIPPS observation manual.
%
%  TO ADD OR EDIT A CONCEPT, copy the template block near the end of this
%  file. A concept page is four calls: \conceptheader, \twopane,
%  \specialnote, and \scoregrid (which takes four arguments).
% =====================================================================
\documentclass[10pt,a4paper]{article}

\usepackage[T1]{fontenc}
\usepackage[utf8]{inputenc}
\IfFileExists{lmodern.sty}{\usepackage{lmodern}}{}
\usepackage[margin=15mm,top=13mm,bottom=13mm]{geometry}
\usepackage[table]{xcolor}
\usepackage{tabularx}
\usepackage{array}
\usepackage{enumitem}
\usepackage{fancyhdr}
\usepackage{ragged2e}
\usepackage[hidelinks]{hyperref}

\usepackage{soul}
\sethlcolor{yellow}

\usepackage{tikz}
\usetikzlibrary{fadings}

% ---------------------------------------------------------------- palette
\definecolor{tnavy}{HTML}{1F3864}
\definecolor{thead}{HTML}{D9E2F3}   % column A/B stem row
\definecolor{tsub}{HTML}{EDF2FA}    % Very / Somewhat row
\definecolor{texfill}{HTML}{F7F7F7} % examples row
\definecolor{trule}{HTML}{9DA9B8}
% Stem-row gradient. Runs warm -> neutral -> cool so the hue itself changes,
% not merely the shade, and so the ordinal direction of the scale is legible
% at a glance. All three stops are light enough for black text on top.
\definecolor{gradA}{HTML}{F7D9BC}   % Column A end: warm sand
\definecolor{gradM}{HTML}{E6E6E2}   % midpoint: neutral
\definecolor{gradB}{HTML}{A8C6E8}   % Column B end: cool blue
\definecolor{tgrey}{HTML}{5A5A5A}

\setlist[itemize]{leftmargin=3.6mm,itemsep=0.9pt,topsep=1.2pt,parsep=0pt,
                  label=\textcolor{tnavy}{\footnotesize$\bullet$}}
 
% Examples list. NOTE: the font size is NOT set with enumitem's before/after
% keys. Doing so leaves the final item's paragraph to be closed while
% \normalsize is active, which gives the last bullet in every cell a larger
% baseline skip than the ones above it. The size is set by \exlistsize inside
% the cell instead, so every item is broken at the same size.
\newlist{exlist}{itemize}{1}
\setlist[exlist]{leftmargin=3.2mm,itemsep=1.6pt,topsep=0pt,parsep=0pt,
                 label=\textcolor{tgrey}{$\bullet$}}
\newcommand{\exlistsize}{\scriptsize}

\pagestyle{fancy}\fancyhf{}
\renewcommand{\headrulewidth}{0pt}
\fancyfoot[C]{\footnotesize\color{tgrey}\thepage}
\setlength{\parindent}{0pt}
\setlength{\tabcolsep}{2.2mm}
\renewcommand{\arraystretch}{1.25}

% ---------------------------------------------------- concept page header
% Large numeral to the left of the concept statement, as in TIPPS.
\newcommand{\conceptheader}[2]{%
  \noindent
  \begin{minipage}[c]{13mm}\raggedright
    {\fontsize{30}{30}\selectfont\bfseries\color{tnavy}#1}
  \end{minipage}%
  \begin{minipage}[c]{\dimexpr\linewidth-13mm\relax}\RaggedRight
    {\large\bfseries #2}
  \end{minipage}\\[1.5mm]
  \textcolor{tnavy}{\rule{\linewidth}{1.1pt}}\\[2mm]}

% ------------------------------------------- importance | indicators pair
\newcommand{\twopane}[2]{%
  \noindent
  \begin{minipage}[t]{0.475\linewidth}\RaggedRight
    {\bfseries\small Importance of Concept}\\[1.4mm]
    \footnotesize #1 \normalsize
  \end{minipage}\hfill
  \begin{minipage}[t]{0.475\linewidth}\RaggedRight
    {\bfseries\small Indicators include:}\\[0.6mm]
    #2
  \end{minipage}\\[2mm]}

% ------------------------------------------------------------ special note
\newcommand{\specialnote}[1]{%
  \noindent\textcolor{trule}{\rule{\linewidth}{0.4pt}}\\[1.2mm]
  {\footnotesize\textbf{SPECIAL NOTE:} #1}\\[2mm]}

% ------------------------------------------------------- the scoring table
% \scoregrid{Column A stem}{Column B stem}{anchor row}{examples row}
% The two row arguments are four ampersand-separated cells each.
%
% STEM ROW: drawn as ONE full-width tikzpicture inside a \multicolumn{4} cell
% with @{} on both sides, so the fill reaches the cell edges exactly and no
% white stripe can appear. A horizontal gradient runs from gradA (Column A
% end) to gradB (Column B end), which also encodes the direction of the scale.
% Height is measured from the two text boxes so the band always fits.
\newlength{\gridw}\newlength{\stemw}\newlength{\stemh}\newlength{\stemtop}

% The array package inserts \@arstrut at the start of every row: a strut about
% \arraystretch x \ht\strutbox tall ABOVE the baseline. A tikzpicture whose
% bounding box starts at y=0 therefore leaves that strut unpainted, which shows
% as a white band across the top of the cell. The fix is to extend the shaded
% rectangle upward by \stemtop so the fill covers the strut as well as the
% text, and to hang the text from the top edge rather than from the baseline.
\newsavebox{\stemAbox}\newsavebox{\stemBbox}

\newcommand{\stemrow}[2]{%
  \setlength{\stemtop}{1.15\ht\strutbox}%
  \setlength{\stemw}{0.5\dimexpr\gridw-2\arrayrulewidth\relax}%
  \sbox{\stemAbox}{\begin{minipage}[t]{\dimexpr\stemw-4.4mm\relax}\RaggedRight
      \footnotesize\textbf{Column A}\\#1\end{minipage}}%
  \sbox{\stemBbox}{\begin{minipage}[t]{\dimexpr\stemw-4.4mm\relax}\RaggedRight
      \footnotesize\textbf{Column B}\\#2\end{minipage}}%
  \setlength{\stemh}{\dimexpr\ht\stemAbox+\dp\stemAbox\relax}%
  \ifdim\dimexpr\ht\stemBbox+\dp\stemBbox\relax>\stemh
    \setlength{\stemh}{\dimexpr\ht\stemBbox+\dp\stemBbox\relax}\fi
  % total band = strut above baseline + top pad + text + bottom pad
  \setlength{\stemh}{\dimexpr\stemtop+\stemh+2.6mm\relax}%
  \begin{tikzpicture}[inner sep=0pt,outer sep=0pt,baseline=0pt]
    \shade[left color=gradA,middle color=gradM,right color=gradB]
      (0,\stemtop) rectangle (\dimexpr\gridw-2\arrayrulewidth\relax,
                              \dimexpr\stemtop-\stemh\relax);
    \draw[black,line width=\arrayrulewidth]
      (\stemw,\stemtop) -- (\stemw,\dimexpr\stemtop-\stemh\relax);
    \node[anchor=north west] at (2.2mm,\dimexpr\stemtop-1.2mm\relax)
      {\usebox{\stemAbox}};
    \node[anchor=north west] at (\dimexpr\stemw+2.2mm\relax,
                                 \dimexpr\stemtop-1.2mm\relax)
      {\usebox{\stemBbox}};
  \end{tikzpicture}}

\newcommand{\scoregrid}[4]{%
  \setlength{\gridw}{\linewidth}%
  \renewcommand{\arraystretch}{1.06}%
  \noindent\footnotesize
  \begin{tabularx}{\linewidth}{|>{\RaggedRight\arraybackslash}X
                               |>{\RaggedRight\arraybackslash}X
                               |>{\RaggedRight\arraybackslash}X
                               |>{\RaggedRight\arraybackslash}X|}
  \hline
  \multicolumn{4}{|@{}l@{}|}{\stemrow{#1}{#2}}\\
  \hline
  \rowcolor{tsub}
  \centering\arraybackslash\textbf{Very Accurate} &
  \centering\arraybackslash\textbf{Somewhat Accurate} &
  \centering\arraybackslash\textbf{Somewhat Accurate} &
  \centering\arraybackslash\textbf{Very Accurate}\\
  \hline
  #3\\
  \hline
  \rowcolor{texfill} #4\\
  \hline
  \end{tabularx}\normalsize}

\begin{document}


\thispagestyle{empty}
\vspace*{18mm}
\begin{center}
{\Huge\bfseries\color{tnavy} LTL}\\[5mm]
{\LARGE\bfseries Observation Manual for Secondary Classrooms}\\[4mm]
{\large Uganda --- Summarized TIPPS-based instrument}\\[14mm]

\textcolor{trule}{\rule{0.62\linewidth}{0.4pt}}\\[8mm]

\end{center}
\vfill
\begin{center}
\footnotesize
This instrument is adapted from the Teacher Instructional Practices and Processes System (TIPPS).\\[2mm]
Seidman, E., Kim, S., Raza, M., Ishihara, M., \& Halpin, P.~F. (2018). Assessment of pedagogical practices
and processes\\ in low and middle income countries: Findings from secondary school classrooms in Uganda.
\textit{Teaching and\\ Teacher Education}, 71, 283--296. \texttt{doi:10.1016/j.tate.2017.12.017}\\[1.5mm]
Seidman, E., Raza, M., Kim, S., \& McCoy, J.~M. (2013). \textit{Teacher Instructional Practices and\\
Processes System (TIPPS): Manual and scoring system.} New York University.
\end{center}
\clearpage


\conceptheader{A}{Observer Manual}
\footnotesize
This instrument observes eight concepts of teacher practice and classroom process. For each concept the
manual provides three things. First, an \textbf{Importance of Concept} paragraph: understanding the
concept's deepest meaning is what makes the choice possible. Second, \textbf{indicators} that can be
readily observed in the classroom --- this is not an exhaustive list, and an observer is not expected to
find all of the indicators in any one observation period. Third, \textbf{examples for all four options},
illustrating the extent to which the practice may or may not be occurring. The examples given may or may
not occur in the classroom being observed and should not be relied on as the only manifestations of the
practice, but as possible ones.

\vspace{2mm}
Approach each concept from the top down, beginning with the importance of the concept rather than moving
directly to the examples. Take detailed notes throughout the observation. Then, for each concept, review
your notes against the concept and its indicators and decide:

\vspace{2mm}
\noindent\textbf{Step 1.} Which column, A or B, best reflects the classroom? Think about the concept and
the indicators. If the concept were split down the middle, are its characteristics more or less present?

\vspace{1mm}
\noindent\textbf{Step 2.} Within your chosen column, is it Somewhat Accurate or Very Accurate? Column A
Very Accurate generally represents the absence of the concept; Column B Very Accurate represents its
presence in the form that most closely resembles the indicators and the importance of the concept itself.
Match your notes against the anchor description in each cell, and use the examples only to calibrate.
\normalsize

\vspace{4mm}
\noindent\textcolor{tnavy}{\rule{\linewidth}{0.8pt}}\\[2mm]
{\bfseries\small Guiding rules for coding}\\[1.5mm]

\vspace{1.1mm}\noindent{\footnotesize\textbf{The scale.}\ Four options, lowest to highest: Column A Very Accurate (the practice is clearly absent), Column A Somewhat Accurate, Column B Somewhat Accurate, Column B Very Accurate (the practice is strongly present). Column B is always the positive end. Decide the column first, then the option within it.}\\

\vspace{1.1mm}\noindent{\footnotesize\textbf{Start at Column A, Very Accurate.}\ Assume the practice is not occurring and move off that assumption only for something you actually saw or heard. Where the evidence is missing, ambiguous, or off-camera, the score is Column A. Do not fill the gap with an assumption.}\\

\vspace{1.1mm}\noindent{\footnotesize\textbf{Evidence for any Column B score.}\ Every Column B score needs one concrete moment in your notes: the time, and what was said or done. No moment, no Column B. A Very Accurate score in Column B needs the practice to recur across the lesson, never a single moment however strong.}\\

\vspace{1.1mm}\noindent{\footnotesize\textbf{Read the Special Note first.}\ Each item's Special Note lists what does not count towards it. If the lesson offers only things on that list, the score stays in Column A regardless of how good the rest of the lesson looks.}\\

\vspace{1.1mm}\noindent{\footnotesize\textbf{Score each concept on its own.}\ Do not look at your other scores, and do not let a strong impression of the teacher carry across. A warm, orderly, energetic class can be Column A on every concept in this manual. Where two concepts look similar --- concepts 2 and 3 especially --- score concept 3 first, and score it on what the pupils did regardless of how the teacher responded.}\\

\vspace{1.1mm}\noindent{\footnotesize\textbf{Look for what would lower the score.}\ Before settling on Column B, name what would have made it Column A: the teacher stepping out, the answer being supplied, the check never happening. Write it in your notes even when you still score Column B.}\\

\vspace{1.1mm}\noindent{\footnotesize\textbf{Indicators are not a checklist.}\ The indicators are examples of what the practice looks like, not a list to tick off. This is not an exhaustive list, and an observer is not expected to find all of the indicators in an observation period. A lesson may show none of them and still earn the option; a lesson may show several and still miss it. Judge the concept.}\\

\vspace{1.1mm}\noindent{\footnotesize\textbf{Examples are illustrative only.}\ The examples in the scoring table are possible manifestations of an option, not templates. Do not require a lesson to match an example, and do not penalise it for looking different.}\\

\vspace{1.1mm}\noindent{\footnotesize\textbf{Language and recording.}\ Audio may move between English, Luganda, and Lusoga, and the transcript may be garbled. Never mark a teacher down for speech-recognition errors, but do not credit a practice because the record is unclear either. The teacher has left the room only if you see them go out through the door: walking to the back of the room, or out of frame, is still present. Brief absences of a minute or two do not count as leaving.}


\clearpage
\conceptheader{B}{How much of the class --- the shared reach scale}
{\footnotesize One reach scale is used throughout this manual. Where an item refers to how much of the class was involved, it uses one of these four phrases and nothing else. Judge reach by counting distinct pupils in your notes as you watch and converting at the end: a count can be checked afterwards, an impression cannot. A group the teacher genuinely worked with counts as its members reached.}\\[2mm]

\vspace{0.9mm}\noindent{\footnotesize\textbf{Almost no one.}\ One or two pupils while the rest of the class watches, or the same single pupil returned to repeatedly.}\\

\vspace{0.9mm}\noindent{\footnotesize\textbf{A few of the same pupils.}\ A small group carries it throughout --- the same voices, or the same corner of the room. You could name them.}\\

\vspace{0.9mm}\noindent{\footnotesize\textbf{A good share of the class.}\ Clearly beyond the same few. Pupils from different parts of the room take part, though still well short of everyone. In a class of 60, roughly 15 pupils or more.}\\

\vspace{0.9mm}\noindent{\footnotesize\textbf{Most of the class.}\ You would struggle to point to many pupils who were left out. Includes methods designed to cover everyone --- all pupils writing, all groups visited, most participating.}

\vspace{3mm}
\noindent{\scriptsize\color{tgrey}\textit{Decision record: this scale replaces the percentage thresholds
used in the detailed AI rubric. Those thresholds all sat in the same region and nothing in the concepts
justified different cut-points, so one scale is now used identically throughout. Verbal quantifiers are
read differently by different people, so the bands must be calibrated against master-coded lessons before
live coding begins, and the AI rubric must be moved to the same scale.}}
\clearpage

\begin{center}\vspace*{60mm}{\Huge\bfseries\color{tnavy}Concepts}\end{center}
\clearpage



% ===================================================== BEGIN CONCEPT 1 ===
\conceptheader{1}{Teacher creates opportunities for cooperative and collective learning.}
 
\twopane{%
Learning together means that pupils and teacher are in genuine conversation about the substance of the
lesson. Two arrangements qualify equally. Pupils may work with and from one another in pairs or groups,
explaining, comparing, and building on each other's thinking. Or the class and the teacher may hold a
sustained exchange in which many different pupils contribute and the teacher works with what they offer
rather than accepting it and moving on. Neither arrangement is worth much on its own. What makes
collaboration productive is the teacher facilitating it: telling pupils how to work with one another,
modelling what a good explanation, question, or response looks like, moving among them while they work,
and drawing their contributions together. A heavily teacher-directed lesson can therefore still show this practice at its fullest, provided the teacher is in real conversation with pupils and the questions build on one
another. 
}{%
\footnotesize
\begin{itemize}
  \item Teacher tells pupils how to work with one another --- compare your answers, explain your method to
        your partner, check each other's work --- and not only what task to do
    \item Teacher demonstrates the exchange itself before pupils attempt it, modeling how to question a
        partner or respond to a classmate
  \item Teacher moves among the pairs or groups while they work, leaning in, questioning, and looking at
        what they are producing
  \item Pupils are prompted to build on, respond to, or give feedback on a classmate's answer
  \item Teacher takes a pupil's contribution in whole-class exchange and works with it, probing it, linking
        it to another pupil's, or turning it back to the class
  \item Teacher questions, or opens discussion around, a pupil who is presenting or working at the board
\end{itemize}
\normalsize
}
 
\specialnote{Choral reading and call-and-response, pupils seated in groups but working through the task
individually, and a pupil presenting or writing at the board with no exchange around it do not count as
cooperative or collective learning, however much of the lesson they occupy.}
 
\scoregrid{Teacher does not create opportunities for cooperative or collective learning.}%
          {Teacher creates opportunities for cooperative or collective learning.}{%
% ---- anchor descriptions, left to right ----
\footnotesize No opportunity is created. The lesson runs on lecture, individual work, choral answering, or
one pupil at a time answering the teacher. Where a pupil presents or works at the board, the teacher asks
nothing and opens no discussion around it. &
\footnotesize An opportunity is mentioned but not built: no clear objective, no real time given, or only a
few of the same pupils carry every exchange while the rest stay passive, so that fewer than a good share of the
class ever gets a chance to contribute. &
\footnotesize A real exchange occurs and a good share of the class is drawn in, but the teacher sets it up
and then withdraws: little guidance on how to work together, little follow-up on what pupils produce, little moving
among the groups, only one group engaged, or the teacher out of the room for much of it without a purpose
that serves the activity.  &
\footnotesize A real exchange, a good share of the class involved, and the teacher facilitating throughout.
Pupils are told, and shown where needed, how to work with one another; the teacher reaches more than one
group at more than one point; and contributions are built on rather than accepted and passed over.
}{%
% ---- examples, left to right ----
\exlistsize\begin{exlist}
  \item The lesson runs on choral reading and call-and-response drills throughout.
  \item A pupil is sent to the board to write an answer. The teacher asks nothing, invites no response from
        the class, and calls the next pupil.
  \item Pupils work through an exercise individually while the teacher marks at the desk.
\end{exlist} &
\exlistsize\begin{exlist}
  \item The teacher tells pupils to discuss a question with the person next to them, allows a few seconds,
        then announces the answer himself.
  \item Each group reads out its answer in turn. The teacher thanks each one and calls the next, never
        asking the class to compare or question them.
  \item A discussion runs well, but the same three pupils supply every contribution.
\end{exlist} &
\exlistsize\begin{exlist}
  \item Pupils work together as told, but the teacher stays with one group throughout and never visits the
        others.
  \item Many pupils contribute to a whole-class exchange, but each answer is accepted and passed over.
  \item The teacher sets the groups working and leaves the room for a significant part of the activity, returning only to
        collect answers.
\end{exlist} &
\exlistsize\begin{exlist}
  \item ``Compare your answers, then explain your method to your partner'' --- the teacher moves between
        groups, leans in, and later weaves pupils' answers into the discussion.
  \item The teacher demonstrates how to question a partner's method, then visits several pairs while they
        work.
  \item In discussion the teacher asks one pupil to respond to another's answer, then asks a third whether
        they agree and why.
\end{exlist}
}
% ======================================================= END CONCEPT 1 ===
 
\clearpage

% ===================================================== BEGIN CONCEPT 2 ===
\conceptheader{2}{Teacher creates opportunities for pupils to exercise intellectual agency, and pupils
take them up.}
 
\twopane{%

A teacher that puts intellectual decisions in pupil's hands gives 
This concept looks at what the teacher does to put real intellectual decisions in pupils' hands, and at
whether those attempts work. A teacher creating agency sets questions and tasks that have more than one
defensible answer or approach, gives pupils enough time to think about them without rushing, and then draws
pupils in: inviting them to share the approach they took, asking others to respond to it, probing for more
rather than supplying the answer. Asking pupils to supply examples of a category the teacher has already
defined, or to apply the procedure just demonstrated, is not a real intellectual decision, however willingly
pupils comply. Creating the opportunity is not enough on its own. The teacher must actively promote it, and
pupils must actually take it up.  An attempt that falls flat, or one extended to only one or two pupils while
the rest watch, has not yet become this practice; how much of the class is enabled to exercise agency is what
distinguishes its fullest form.%
}{%
\footnotesize
\begin{itemize}
  \item Teacher sets a question or task with more than one defensible answer or approach, and allows time to
        think before taking answers
  \item Teacher invites pupils to share the approach they took, and then works with what they offer
  \item Teacher asks other pupils to comment on, add to, or evaluate a classmate's answer
  \item Teacher prompts and rephrases without supplying the answer when the class goes quiet
  \item Teacher follows a pupil's question where it leads instead of returning to a fixed agenda
  \item Teacher looks at what pupils produced during the activity rather than setting it and moving on
\end{itemize}
\normalsize
}
 
\specialnote{Choral reading and unison answering are never evidence of agency, however energetic, but in
moderation they are not a fault; they place a lesson in Column A only when they dominate it. Where the
evidence falls between the two Somewhat Accurate options, score Column A.}
 
\scoregrid{Teacher does not create opportunities for pupils to exercise intellectual agency, or the
opportunities do not succeed.}%
          {Teacher creates opportunities for pupils to exercise intellectual agency and pupils take them
up.}{%
% ---- anchor descriptions, left to right ----
\footnotesize No opportunity is created. The task is fully teacher-directed, with one expected method or
response and no space for pupils to generate ideas, or the lesson is dominated by choral reading and unison
answering with little time anywhere for independent thinking. &
\footnotesize An opportunity is offered but does not succeed. The teacher asks for ideas or sets pupils to
try different approaches, but gives no thinking time, offers no support while they work, keeps the task too
constrained for real exploration, or extends the opportunity to only one or two pupils while the rest watch.
Opportunities that only ask pupils to supply examples of a category the teacher has defined, or to apply the
procedure just taught, also belong here. &
\footnotesize Agency opportunities are clear and sustained, the teacher promotes them by encouraging pupils
to share their thinking and by asking questions and making comments that push them further, and pupils take
them up --- but fewer than a good share of the class is enabled to exercise agency. It makes no difference
whether the opportunities are concentrated in one phase of the lesson or spread through it. &
\footnotesize Agency opportunities are clear and sustained, and a good share of the class is enabled to
exercise agency. The teacher encourages pupils to share their thinking and supports that sharing, asking
questions and making comments that push them further, and brings pupils into responding to and evaluating
one another's answers.
}{%
% ---- examples, left to right ----
\exlistsize\begin{exlist}
  \item Every question has one expected answer, and the lesson runs on group reading and choral response
        from start to finish.
  \item The teacher asks only closed questions and answers several of them himself.
  \item A pupil copies answers from a book onto the board. Neither the teacher nor the class engages with
        what he wrote.
\end{exlist} &
\exlistsize\begin{exlist}
  \item The teacher gives some background and then asks for examples of the topic. Pupils supply them and
        the lesson moves on.
  \item Pupils copy the teacher's procedure step by step and never work a problem for themselves.
  \item The teacher tells pupils to work together but is absent while they do. He later asks whether anyone
        took a different approach, no one answers, and he moves on without looking at what they wrote.
\end{exlist} &
\exlistsize\begin{exlist}
  \item The teacher sets a problem on the board and encourages pupils to attempt it, supporting each one who
        comes up, but only a handful come up across the whole lesson.
  \item A whole-class discussion draws genuine contributions with the teacher's encouragement, but from well
        under a good share of the class.
  \item The teacher asks pupils to suggest their own methods. Two do so with his support while the rest copy
        the result.
\end{exlist} &
\exlistsize\begin{exlist}
  \item ``What are the challenges of work today?'' The teacher waits, prompts by referring to pupils' own
        communities without supplying answers, asks ``how?'' when a pupil says ``insecurity'', and different
        pupils build on each other's points.
  \item Several pupils attempt parts of a problem at the board with the teacher's encouragement; he then
        sets group work and moves among the groups, so that between the two a good share of the class takes
        part.
  \item After a pupil presents, the teacher asks the class what they would question about it, and several
        pupils do.
\end{exlist}
}
% ======================================================= END CONCEPT 2 ===
 

\clearpage


% ===================================================== BEGIN CONCEPT 3 ===
\conceptheader{3}{Pupils independently exercise intellectual agency.}
 
\twopane{
Intellectual initiative taken by pupils themselves shows what they are willing and able to contribute when nothing obliges them to --- an idea, a question, an alternative method or representation, a connection, or a challenge to what has been said. The pupil's action counts even when the teacher ignores it, rejects it, or moves past it. Initiative of this kind can appear in many forms; the three seen most often  are a single pupil taking initiative, pupils exchanging and testing ideas with one another, and a class eager to take part when an opportunity opens. \hl{Any one of them is enough, and initiative that resembles none of  them still counts, provided it is the pupil's own contribution and is relevant and thoughtful. Throughout,  what matters is who started it.} A pupil who raises a hand unasked, or who brings up a point nobody invited, is taking initiative. A pupil who answers because the teacher called their name has done what was asked of them, which is not evidence for this concept.
}{
\footnotesize
\begin{itemize}
  \item A pupil asks a substantive question that nobody invited
  \item A pupil offers an interpretation, method, or solution different from the one being presented
  \item A pupil challenges or questions a point the teacher or another pupil has made
  \item Pupils build on, question, or disagree with each other's ideas and give reasons for doing so
  \item Many hands go up, or several pupils volunteer for the board, when an opportunity opens
  \item Pupils pursue an implication beyond what was asked of them
\end{itemize}
\normalsize
}
 
\specialnote{Correct answers, rewording an expected answer, and choral responses are not initiative. Noise
and animation are not peer dialogue: you must be able to point to a pupil engaging with another pupil's
idea. The same few confident pupils answering repeatedly is not eagerness.}
 
\scoregrid{Pupils do not independently exercise intellectual agency.}%
          {Pupils independently exercise intellectual agency.}{%
% ---- anchor descriptions, left to right ----
\footnotesize No pupil-generated intellectual contribution is observed. Pupils recall, repeat, copy, and
follow a teacher-supplied procedure, answering closed questions with short answers. Peer interaction is
absent, or limited to logistics, sharing materials, or working independently while seated in groups. &
\footnotesize An isolated or weak instance of initiative occurs. A pupil volunteers an unprompted comment or
question, but it is minor and procedural rather than about the ideas in the lesson. &
\footnotesize At least one clear instance of independent intellectual agency occurs. A pupil introduces a
substantive unprompted idea, question, solution, or extension, or one of the routes shows a moderate degree
of agency, though the evidence is isolated rather than repeated. &
\footnotesize Several clear instances of independent intellectual agency occur, or fewer but strong ones:
unprompted substantive questions, or pupils suggesting ideas that challenge or question the teacher. It may
also show as pupils engaging with one another's thinking in group work, asking the teacher several
questions, and participating eagerly and offering answers liberally. Pupils quickly volunteering, or getting
up to go to the board to solve a problem when the teacher asks, is also evidence.
}{%
% ---- examples, left to right ----
\exlistsize\begin{exlist}
  \item Pupils answer the teacher's closed questions with short answers and copy the notes from the board or
        from his dictation.
  \item Pupils take turns reading assigned sentences in chorus, asking nothing and offering no unprompted
        comment.
  \item Pupils share equipment, but each completes the worksheet independently using the teacher's
        instructions.
\end{exlist} &
\exlistsize\begin{exlist}
  \item A pupil asks whether the class may use a calculator, \hl{or to go to the bathroom}. % from RUBRIC_and_PROTOCOL_v2 - not sure how much agendy there's in asking to go to the bathroom, I'm not really sure why we had it, since students never ask for permission and just quietly stand up and leave the classroom hmmm
  \item Pupils decide who will write and who will present a group task, but not how to solve or interpret
        it.
  \item A pupil asks the teacher to repeat the last sentence so that they can finish copying it.
\end{exlist} &
\exlistsize\begin{exlist}
  \item The teacher states something and a pupil responds ``But I thought\ldots{}'', questioning the point
        he has made.
  \item A group notices an unusual measurement and asks about it.
  \item After a group activity, several pupils eagerly volunteer to go to the board to present their answers
        once the teacher opens the opportunity.
\end{exlist} &
\exlistsize\begin{exlist}
  \item Reviewing a past paper in biology, pupils freely ask about twins, zygotes, and chromosomes. One
        pushes back --- ``I thought it was the same sperm'' --- and most of the class disagrees, gesturing
        that it takes two.
  \item A pupil asks, unprompted, a substantive question the teacher had not invited; another offers an
        interpretation different from the one being presented, and the direction of the lesson shifts to
        engage with it.
  \item Pupils propose and compare solutions and pursue implications beyond the original prompt, sharing
        ideas liberally throughout the lesson.
\end{exlist}
}
% ======================================================= END CONCEPT 3 ===
 

\clearpage


% ===================================================== BEGIN CONCEPT 4 ===
\conceptheader{4}{Teacher develops critical thinking and deeper learning.}
 
\twopane{%
Deeper learning is happening when the class does genuine intellectual work with the content rather than
restating it. That work most often takes one or a combination of six forms: reasoning a claim through, 
applying an idea to a new context, breaking information down and interpreting it, connecting one concept 
to another so that it explains something, weighing alternatives with reasons, and giving pupils tools to 
structure their own thinking. The claims, methods, and examples being worked on may come from the
teacher or from a pupil, and either may do more of the talking, provided the class is genuinely working
on them together. Two things are needed together: the work must happen on more than one occasion, since an 
isolated episode is not enough, and pupils must be part of it, since a one-way explanation, however clear, 
is not deeper learning. Pupils are part of it when they contribute audibly, share approaches from 
their own understanding, are drawn in to have their thinking tested, or when a pupil's question sparked it. 
Definitions given and examples collected are recall, not depth.%
}{%
\footnotesize
\begin{itemize}
  \item A claim is justified or worked through rather than simply stated --- ``why is that?'', ``how do you
        know?'', ``you said X, but how could that be if Y?''
  \item Something already learned is used on a new example or situation not discussed before
  \item Information is broken into parts and the class discusses what pupils understand from it
  \item One concept is used to explain another, and the class engages with the link
  \item Two answers, methods, or examples are put side by side and compared with reasons
  \item Pupils write their understanding in their own words, structure their approach, or compare
        answers with a partner
  \item A pupil's question triggers an extended explanation, and the teacher checks back with the class
        as it builds
\end{itemize}
\normalsize
}
 
\specialnote{Recall and closed questions do no harm: a lesson with many of them can still be scored in
Column B, provided deeper learning also occurs. But questions alone are not critical thinking, and repeating
a pupil's words is not extending them.}
 
\scoregrid{Teacher does not develop critical thinking and deeper learning.}%
          {Teacher develops critical thinking and deeper learning.}{%
% ---- anchor descriptions, left to right ----
\footnotesize The interaction between teacher and pupils is primarily one-way; the teacher is the only one
speaking. Or the lesson consists of definition-giving and recitation with nothing done with the definitions. &
\footnotesize Deeper-learning prompts occur but are limited, engaging only a few pupils. What the teacher
sets up is not spread through the lesson and amounts to only one or two instances. The teacher does not
involve pupils enough for them to develop their understanding: no audible pupil contributions, no sharing of
approaches, no writing from their own understanding, no checking that pupils are following, and no pupil
question that sparked it. &
\footnotesize The teacher promotes some, but limited, critical thinking. Deeper-learning prompts occur ---
open questions, a request to justify, an application task, a pupil-sparked explanation --- but they are
brief, engage only a few pupils, are not opened to the class, or the reasoning is curtailed by the teacher
supplying the answer. &
\footnotesize Several sustained deeper-learning episodes in which pupils are engaged and the teacher
supports the reasoning through to a built-up understanding. The teacher works to involve a good share of the
class across the lesson, rather than conducting the episode with one or two pupils while the rest watch.
Pupils may contribute audibly, share their thinking, write their own understanding, or compare answers with
a partner.
}{%
% ---- examples, left to right ----
\exlistsize\begin{exlist}
  \item ``Repeat after me: when it is raining I must use an umbrella, otherwise I will get wet.'' Pupils
        repeat it again and again. ``Now write down the full answer,'' says the teacher.
  \item The teacher asks ``what is photosynthesis?'', a pupil recites the definition, the teacher links it
        to the definition of respiration, and moves on.
  \item The teacher's only questions finish his own sentences (tag questions) --- ``the leaf makes its food in the\ldots{} what?'' --- which pupils complete in chorus.
\end{exlist} &
\exlistsize\begin{exlist}
  \item ``When it is raining you must use an umbrella. Do you know why? Because otherwise you will get
        wet.'' The teacher points to a pupil: ``Tell me when I must use an umbrella?'' The pupil answers
        ``when it is raining''. ``I will get wet otherwise --- that is the full answer, write it down.''
  \item ``Who can give me an example of a mammal?'' Five answers are collected and each is affirmed, but
        none is explored or contrasted and the lesson moves on.
  \item Many pupils write answers on the board, but the teacher never probes the reasoning behind
        them or checks that the rest of the class is following.
\end{exlist} &
\exlistsize\begin{exlist}
  \item ``Your mother asks you to take an umbrella. What does that mean?'' A pupil answers ``because it is
        raining''. ``Good,'' says the teacher, and puts the same question to another pupil without taking
        it further.
  \item The teacher asks a genuine ``why?'' and one pupil reasons it through, but the exchange stays between
        the two of them and is never opened to the class.
  \item A pupil's question triggers a genuinely deep explanation, but the teacher never checks back with the
        questioner or draws the rest of the class into it.
\end{exlist} &
\exlistsize\begin{exlist}
  \item ``If your mother asks you to take an umbrella, what does that tell you about the weather?''
        Several pupils respond. One says it is raining; another says it could also be very hot. ``Great
        point --- let us think about the uses of an umbrella,'' says the teacher, drawing more pupils in.
  \item A pupil uses an imprecise word in a comment. The teacher takes it up with the whole class, explains
        a better word the pupil could have used and why it fits, and works through the difference.
  \item ``Is agriculture a business?'' Pupils say no. Instead of correcting them the teacher asks ``do the
        cows take themselves to the garden?'', and has pupils work out how manure reaches the crops until
        the class can explain why agriculture involves work and cost like any business.
\end{exlist}
}
% ======================================================= END CONCEPT 4 ===
 

\clearpage


% ===================================================== BEGIN CONCEPT 5 ===
\conceptheader{5}{Teacher provides scaffolding to support pupil understanding.}
 
\twopane{%
Scaffolding is the support that makes a lesson build on itself. The teacher sets out what is to be learned,
explains the concept or method in language pupils can follow, draws them in to check they are following,
anticipates the points they will find hard, and leaves room for them to ask questions --- and then supports
them while they apply it. Explanation alone is not scaffolding, however fluent and accurate, and neither is
setting a task and withdrawing from it. The lesson should build: the teacher prepares pupils for what they
are about to attempt, pupils then attempt it, and the teacher is present and responsive while they do ---
breaking things into steps, prompting, guiding, targeting the particular point where pupils are stuck, and
changing the approach when the first attempt does not land. A lecture-style opening, including one with
choral responses, is a normal way to begin and is not a fault. What matters is what happens once pupils
start applying.%
}{%
\footnotesize
\begin{itemize}
  \item Teacher explains what is to be learned and how the method works, in language pupils can follow
  \item Teacher breaks a concept, method, or problem into steps that are named and explained
  \item Teacher builds on knowledge pupils already have, saying which earlier work is relevant and how
  \item Teacher offers hints, cues, analogies, or guiding questions while pupils are working
  \item Teacher works initial examples with the class to check pupils can manage parts of the problem
  \item Teacher elaborates on the part pupils find difficult and connects it back to the larger concept
  \item Teacher changes the explanation, example, or representation when pupils remain confused
  \item Teacher reduces assistance so pupils apply, explain, or give each other feedback more independently
\end{itemize}
\normalsize
}
 
\specialnote{A pupil at the board is evidence about this concept only through what the teacher does with it:
working through the error with them, asking questions during the work, or opening it to the class. Watching
and then making a superficial correction is not scaffolding.}
 
\scoregrid{Teacher does not provide scaffolding to support pupil understanding.}%
          {Teacher provides scaffolding to support pupil understanding.}{%
% ---- anchor descriptions, left to right ----
\footnotesize The teacher provides no meaningful scaffold. Content, procedures, or definitions are stated,
demonstrated, or dictated without making the reasoning accessible to all. Observable confusion is ignored,
or the same explanation is repeated unchanged. Important language may be left unexplained. Or the lesson is
a sequence of disconnected activities with no stated purpose and no visible build. &
\footnotesize The teacher may lecture, provide definitions, or ask pupils to supply some definitions and
information, and then, after some more general explanation, give pupils a problem to solve. If the teacher
does not support pupils through the problem-solving phase, it is not scaffolding. &
\footnotesize The teacher provides clear and evident support to prepare pupils to solve a problem or
complete an activity, and works through examples as a class that are quite teacher-driven --- but then does
not allow them much space to try for themselves or to ask questions. The teacher checks whether pupils can
solve parts of a problem, for example by running through initial examples on the board together. &
\footnotesize The teacher is present and around, so support is readily available. He gives clear and evident
support to prepare pupils, then allows them space to try and to ask questions, and supports them through
that process, adjusting the approach if it does not land. Where the sequence permits, assistance is reduced
so pupils apply, explain, or give each other feedback more independently; equipping pupils to support one
another counts as scaffolding.
}{%
% ---- examples, left to right ----
\exlistsize\begin{exlist}
  \item The teacher demonstrates double-digit multiplication on the board but never explains the steps or
        gives pupils a chance to work through them.
  \item The teacher dictates the definitions of ``equilateral'' and ``right-angled'' triangle without
        explaining or illustrating the terms.
  \item Pupils show they are confused, but the teacher repeats exactly the same explanation and continues.
\end{exlist} &
\exlistsize\begin{exlist}
  \item The teacher demonstrates double-digit multiplication with little elaboration. When pupils ask
        questions he points at the example and says ``do that step correctly'', or ``look at how it is
        broken down on the board''.
  \item ``Think about what we did yesterday,'' the teacher tells a struggling pupil, without saying which
        part of the earlier lesson is relevant or how it helps here.
  \item The teacher gives a fluent, accurate lecture, sets a problem, and is then absent or sitting
        disengaged for most of the time pupils spend working on it.
\end{exlist} &
\exlistsize\begin{exlist}
  \item The teacher demonstrates double-digit multiplication. ``I have broken it into three steps, and each
        must be completed in order.'' When pupils have questions, he repeats the same explanation of all
        three steps.
  \item A pupil solves at the board and the teacher works through the error with them substantively at the
        end --- but asked no questions during the work and did not open it to the class.
  \item The teacher connects the concept to a previous lesson and uses guiding questions, but gives pupils
        no opportunity to apply the connection themselves.
\end{exlist} &
\exlistsize\begin{exlist}
  \item The teacher demonstrates double-digit multiplication. ``I have broken it into three steps.'' When
        pupils have questions he focuses on the step each one needs, reminds the class of problem areas to
        watch out for, and encourages pupils to help each other.
  \item ``Here is a spot where some pupils have difficulty --- let us go through it,'' says the teacher, and
        then goes through it.
  \item The teacher calls a pupil to the board, asks ``what are you missing here?'' while they work, guides
        the revision, then invites the class to comment on each other's solutions.
\end{exlist}
}
% ======================================================= END CONCEPT 5 ===
 

\clearpage


% ===================================================== BEGIN CONCEPT 6 ===
\conceptheader{6}{Teacher checks for understanding.}
 
\twopane{%
Checking for understanding means using methods that actually reveal what pupils have
understood. The test for any method is: did it provide information about the level of understanding of pupils? An action a 
pupil could carry out without understanding --- chanting, copying, repeating a definition --- reveals nothing and does not count towards this concept. A lively class may give a sense of engagement, but participation alone is not enough. Two things are needed together: the methods must reveal what specific pupils know rather than leaving a general impression, and they must reach a substantial part of the class rather than only the two or three individuals who happen to answer. Checking should also happen at more than one point in the lesson, not in a single isolated moment. Strong evidence comes from the teacher who checks, finds a gap, an error, or a confusion, and acts on it --- clarifying, explaining it differently, asking further questions, turning it into a teaching moment.%
}{%
\footnotesize
\begin{itemize}
  \item Teacher asks pupils to explain a concept in their own terms, or to say why they did what they did
  \item Teacher circulates during application work and engages --- leaning in, questioning, correcting,
        looking at what pupils have written
  \item Teacher has pupils write down what they have understood in their own words, and reviews it
  \item Teacher inspects worked problems as pupils apply a method, rather than while they copy
  \item Teacher sends several pupils to the board to work independently, and questions them while they are there
  \item Teacher questions a group's spokesperson, or has other pupils question them, and gets an answer
  \item Teacher re-explains, rephrases the question, or changes the example on seeing that pupils are not following
\end{itemize}
\normalsize
}
 
\specialnote{Choral and call-and-response work, notebooks inspected while pupils copy from the board, answers written or presented at the board without explanation, and definitions repeated do not count as, checking, however many pupils take part. A lively class is not a check: the teacher must learn what particular pupils know.}
 
\scoregrid{Teacher does not check for understanding.}%
          {Teacher checks for understanding.}{%
% ---- anchor descriptions, left to right ----
\footnotesize The teacher does not ask any questions, give any prompts, or use any strategies to find out
whether pupils understand what is being taught. He lectures or assigns work with no attempt to gauge
comprehension. Any question asked is rhetorical, or takes a choral yes or no that is accepted without
further checking. Call-and-response drilling, however energetic, belongs here when it is the only mode. &
\footnotesize The teacher makes minimal attempts to check understanding, but the methods yield information
about a small section of the class, or he accepts choral and surface-level responses without probing further.
The checking is perfunctory rather than diagnostic. Notebooks inspected while pupils copy from the board or dication,
answers written at the board without explanation, definitions repeated, and a teacher who circulates without
engaging or leaves the class for extended periods all belong here. &
\footnotesize The teacher uses questions, prompts, or other strategies that are effective at determining \hl{a % a bigger part of the classroom? more then half? what's a good number?!?!?!?!
few/some pupils'} level of understanding, but not most. He calls on individuals or asks for volunteers, gaining
information about those who respond, while most pupils' understanding remains unknown. A few \hl{independent/confident} % by confident I mean the ones that are able to present/go to the board without materials to read from
pupils presenting at the front, genuinely explaining what they did, is one of the minimum forms of this
score. Group work left mostly unmonitored with only some groups later checked, also sits here. & 
\footnotesize The teacher's methods, taken together, reach most of the class, counting checked groups as
covered: all pupils writing their understanding and the teacher reviewing it; held-up work scanned;
circulation during application engaging several pupils or groups; spokespersons presenting, combined with
the teacher's questioning and his circulation during the work. In the strongest form, the teacher finds a
gap and acts on it, clarifying, explaining it differently, or teaching into it.
}{%
% ---- examples, left to right ----
\exlistsize\begin{exlist}
  \item The teacher explains how to identify types of angles, asks ``do you all understand?'', the class
        answers ``yes'' in unison, and he moves on without asking anyone to identify one.
  \item The teacher lectures and sets written work, asking nothing about whether pupils have followed.
  \item The class chants sentences in unison for several minutes, and no other method is used all lesson.
\end{exlist} &
\exlistsize\begin{exlist}
  \item After explaining angles the teacher asks ``who can tell me what type of angle this is?'' One pupil
        raises a hand and answers. ``Good,'' says the teacher, and moves on.
  \item A few confident pupils write answers at the board with little or no explanation while the teacher
        watches.
  \item The teacher walks between the desks while pupils copy notes from the board, glancing at books
        without asking anything.
\end{exlist} &
\exlistsize\begin{exlist}
  \item ``Who can give me an example of a reflex angle, or an obtuse angle?'' A few hands go up and the
        teacher calls on two. Their answers are discussed, but the rest of the class is not checked.
  \item The teacher asks three pupils in sequence to explain a concept at the board, gaining information
        about those three but not about the rest.
  \item Groups work without the teacher monitoring them, and only the presenting groups' understanding is
        ever revealed.
\end{exlist} &
\exlistsize\begin{exlist}
  \item The teacher asks all pupils to draw a reflex angle on their slates and hold them up, scans the room
        to see who has drawn it correctly and who is confused, and questions individuals about their books.
  \item During group work the teacher circulates and questions several groups, then each spokesperson
        presents while he probes, so that the class is checked through its groups.
  \item A check reveals that pupils are confused about the second step. The teacher stops, re-explains with
        a new example, and checks again.
\end{exlist}
}
% ======================================================= END CONCEPT 6 ===
 

\clearpage


% ===================================================== BEGIN CONCEPT 7 ===

\conceptheader{7}{Teacher provides specific feedback to facilitate learning.}
\twopane{Specific feedback allows the teacher to gauge the level of understanding of each pupil as well as the class as a whole. Without it, pupils may be guessing about how they are doing and how their work can be improved. Feedback of this kind responds to whatever a pupil offers --- an answer, an idea, a piece of work, a presentation --- and not only to mistakes. What separates feedback that facilitates learning from feedback that does not is whether it moves the pupil towards understanding. A bare verdict does not: telling a pupil ``yes'', ``no'', or ``correct'' settles that one answer but leaves them nothing to think with, and nothing they could use on a different question tomorrow. Feedback that names the rule or reason behind the answer, or that asks pupils what they did and why before building on what they say, gives them something they can carry forward. On an incorrect answer it redirects pupil thinking to the correct process; on a correct answer it reinforces the thinking that got them there.}{\footnotesize\begin{itemize}
  \item Teacher explains why an answer is wrong, or why it is right, in terms of a rule, method, or reason
  \item Teacher targets the correction to the point where the pupil made a wrong turn, and redirects from there
  \item Teacher asks the pupil what they did and why before responding providing the correct answer or building on what they said
  \item The pupil's thinking is surfaced first and the teacher's input visibly builds on it
  \item Teacher works through a correction with the pupil rather than supplying it
  \item Teacher gives substantive comments after a presentation, or while moving among pupils at work
\end{itemize}\normalsize}
\specialnote{Praise and applause are not specific feedback: they neither add to nor take away from this concept.}
\scoregrid{Teacher does not provide specific feedback to facilitate learning.}{Teacher provides specific feedback to facilitate learning.}{%
\scriptsize No feedback is offered during the lesson, or the feedback given is inappropriate. Score here also when the teacher never elicits pupils' work or answers: with nothing offered, there is nothing to give feedback on. &
\scriptsize Feedback occurs but stays at the level of a verdict: the answer is marked right or wrong and nothing more is added. Frequently leaving pupils' questions unanswered belongs here. &
\scriptsize One clear instance of specific feedback --- \hl{a verdict/response/reaction to a pupil that participated} with the reasoning behind it, or a substantive explanation of where an error came from. The teacher is actively eliciting pupils' work or answers so that there is something to respond to. &
\scriptsize Several clear instances, two or more, of specific feedback, with the teacher actively eliciting pupils' input in varied ways so that there are opportunities to give it. At least two of these are full exchanges: the teacher asks what the pupil did and why, clarifies their understanding of the pupil's work, builds on it --- offering insight, supporting the pupil through the thinking, or drawing other pupils into the conversation --- and the exchange ends in a clarification that is clear and understood, with the class brought along when it is public.
}{%
\begin{exlist}
\item The teacher reads the answers to a quiz aloud, without any explanation, for pupils to check against.
\item ``How many times do I have to tell you it is the wrong answer?''
\item The teacher lectures throughout and never asks pupils to produce work or answers, so nothing is ever offered to respond to.
\end{exlist} &
\begin{exlist}
\item ``No, 5 is the wrong answer'' and the lesson moves on.
\item The teacher walks the rows issuing corrections saying ``full stop here'', ``capital letter'', ``wrong sign'', with no reasoning asked for or given
\end{exlist} &
\begin{exlist}
\item ``The answer is 8, not 5 --- remember, addition comes before subtraction here,'' and the teacher moves on.
\item The teacher reviews the most difficult questions on a quiz with the class, showing where common errors were made.
\end{exlist} &
\begin{exlist}
\item ``No, 5 is still not correct, but you are on the right track. Think about the rules --- does addition come before subtraction?'' and the teacher works it through with the pupil.
\item ``Why did you put the full stop there? Is anything missing?'' The teacher listens, works the correction through, then asks the class what changed and why.
\end{exlist}
}


% ======================================================= END CONCEPT 7 ===

\clearpage

% ===================================================== BEGIN CONCEPT 8 ===

\conceptheader{8}{Teacher connects learning to pupils' everyday life.}
\twopane{A strong connection between the lesson and pupils' real-life experiences helps pupils understand what they are studying and see its relevance beyond their time in the classroom. Connecting a concept to something familiar makes it easier to grasp: a pupil who has seen an idea at work in the market, in the garden, or on the walk to school has something concrete to attach it to. What matters is how far the connection is taken. A concept may simply be mentioned alongside something familiar; it may be stated to be connected; it may be explained and developed by the teacher; or pupils may explore it themselves. These are not equally valuable. The connection has to be built rather than gestured at: naming something that exists locally, or setting a problem among local names and prices, leaves the pupil with nothing new to understand. Connections are also not limited to subject content. A lesson may connect just as genuinely to friendship, to handling stressful events, or to skills pupils use outside school.}{\footnotesize\begin{itemize}
  \item Teacher explains how the concept works in a setting pupils know, with enough local detail that the link is unmistakable
  \item Pupils are asked about their own experience of the topic under discussion, they answer and the class works with their contributions
  \item Teacher relates the lesson to experiences pupils are facing, such as how to deal with stressful events
  \item Teacher has pupils carry out in practice the activity being discussed, such as lifting objects of different weights
  \item Teacher directs pupils to talk with one another about a specific experience of their own
  \item Teacher applies the lesson to a skill pupils use outside school
  \item Teacher uses local community examples instead of, or in addition to, examples from the book
  \item Teacher returns to the connection later in the lesson rather than mentioning it once
  
\end{itemize}\normalsize}
\specialnote{Using local names, prices, or market items as the setting of a problem is not a connection, and neither is naming something that exists locally. A practical activity is not by itself a connection to everyday life unless the link to pupils' own experience is made explicit.}
\scoregrid{Teacher does not connect learning to pupils' everyday life.}{Teacher connects learning to pupils' everyday life.}{%
\scriptsize An object or concept that could relate to pupils' lives is mentioned, but no connection is drawn to their experience. Local names, prices, or market items used as the setting of a problem sit here: familiar props, with no link built. &
\scriptsize A connection is made in passing. Pupils might infer it, but the teacher does not make it explicit or develop it, and does not come back to it. &
\scriptsize A genuine, developed connection is articulated by the teacher: the link between the concept and pupils' lives, community, or family is explained clearly and with local detail that pupils could not miss it --- but pupils are given no opportunity or time to discuss, use, or reflect on it themselves. &
\scriptsize The connection is made and pupils get to work with it: discussing it, giving their own examples from everyday life, or practicing it in an activity the teacher organizes. In the strongest form pupils bring up their own experiences and the class explores how the concept relates to them.
}{%
\begin{exlist}
\item ``A tree is a living organism.'' ``A chemical is a substance.''
\item A mathematics problem uses pupils' names and local prices for pens and counter books.
\end{exlist} &
\begin{exlist}
\item ``Chemicals are produced in a factory.'' ``Trees are found in the forest.''
\item A grammar topic opens with ``you talk to your friends every day'' and is never returned to.
\end{exlist} &
\begin{exlist}
\item To explain work, the teacher uses a boy walking to the borehole with a jerrycan and a local flyover the pupils know, then moves on.
\item Percentages are explained through market bargaining with local detail, but pupils never use the connection.
\end{exlist} &
\begin{exlist}
\item ``The forest you walk through on the way to school --- what was that like?'' and pupils answer.
\item After learning about vegetables, the teacher asks which pupils grow, eat, or sell at home and what else they might do with them.
\end{exlist}
}
\clearpage


% ======================================================= END CONCEPT 8 ===


\conceptheader{N}{Concept statement goes here.}
\twopane{Importance of Concept paragraph. Explain what the practice is, why it matters for pupils, and
what the observer is fundamentally judging. Keep it to roughly eighty to a hundred words, written as
prose. Where reach or depth decides the option, say so here in one sentence.}{%
\footnotesize\begin{itemize}
  \item First observable indicator
  \item Second observable indicator
  \item Third observable indicator
  \item Fourth observable indicator
  \item Fifth observable indicator
\end{itemize}\normalsize}
\specialnote{Use this note for anything the prose and the indicators cannot carry: what does not count
towards the concept, and any caveat that would otherwise be mistaken for evidence.}
\scoregrid{Teacher does not \ldots{} (negative stem).}{Teacher does \ldots{} (positive stem).}{%
\footnotesize Anchor description for Column A, Very Accurate. State what absence looks like. &
\footnotesize Anchor description for Column A, Somewhat Accurate. End with the sentence that separates this from Column B. &
\footnotesize Anchor description for Column B, Somewhat Accurate. End with the sentence that separates this from Very Accurate. &
\footnotesize Anchor description for Column B, Very Accurate. State the fullest form of the practice.
}{%
\begin{exlist}\item First example.\item Second example.\end{exlist} &
\begin{exlist}\item First example.\item Second example.\end{exlist} &
\begin{exlist}\item First example.\item Second example.\end{exlist} &
\begin{exlist}\item First example.\item Second example.\end{exlist}
}

\vspace{3mm}
\noindent{\scriptsize\color{tgrey}\textit{Blank template. Copy this page to add a concept. The four
anchor cells read left to right from lowest to highest, and the row beneath each holds its examples.}}
\clearpage


% ======================================================= BEGIN TEMPLATE ===
%
% CONCEPT STATEMENT
%   A full sentence ending in a full stop, naming the actor and the practice.
%   Teacher-framed for every concept except one about what pupils do.
%   The statement establishes who the actor is, so the Importance paragraph
%   must NOT open by naming them again.
%
\conceptheader{N}{Teacher \ldots{} \emph{(full sentence naming the actor and the practice.}}
 
%
% IMPORTANCE OF CONCEPT  --- roughly 150 to 190 words, prose, no bullets.
%   OPEN with the practice, or its effect, as the grammatical subject:
%     "Specific feedback allows the teacher to gauge\ldots"
%     "A strong connection between the lesson and pupils' lives helps\ldots"
%   NOT "The teacher does X when\ldots" --- the statement above said that already.
%
%   CARRIES: what the practice is; why it matters for pupils, stated causally;
%     the decisive test the coder applies, in the source's own question form;
%     the traps (what feels like the practice but is not);
%     a plain-language definition of any term the anchors will reuse;
%     what the fullest form of the practice looks like, in principle.
%
%   NEVER CARRIES: enumerated behaviours (those are indicators);
%     examples (those go in the table);
%     scoring machinery --- "Column B", "the option", "sets the level",
%       "for the practice to be present", "reach the top of";
%     rationale not traceable to the source guidance --- invented purpose
%       clauses read like construct and have no authority behind them;
%     gendered pronouns --- write "the teacher", or use the passive.
%
%   IF TWO CONDITIONS BOTH HOLD, say so as a pair ("two things are needed
%   together: X, and Y"), never in sequence, which reads as the second
%   correcting the first.
%
%   IF THE ORIGINAL TIPPS CARRIES THIS CONCEPT, start from its paragraph.
%
\twopane{%
Open with the practice or its effect as the subject of the sentence, and say plainly what it does for pupils.
Then give the construct: what the practice actually is, and the test the coder applies to decide whether it
happened, phrased as a direct question in the source's own words. Name the traps --- the things that look
like the practice but are not --- and explain, in plain words, any term the anchors below will reuse, so that
a coder reading only the table still understands it. Where two conditions are both required, say that they
are needed together rather than listing them one after the other. Close on what the fullest form of the
practice looks like as a matter of principle, not as a scoring instruction. Keep the whole paragraph between
roughly 150 and 190 words: 190 is what the pane holds. Use no bullets, no examples, no gendered pronouns, and
no words belonging to the scoring apparatus.%
}{%
%
% INDICATORS  --- six to eight items, no terminal punctuation.
%   Begin every item with the concept's own grammatical subject, normally
%     "Teacher \ldots"; a pupil-framed concept begins every item "Pupils \ldots".
%   Each item names a behaviour a coder could point at on the recording.
%   No two items may restate each other.
%   No item may imply a condition the anchors do not impose (listing a
%     combination suggests both parts are required).
%   No item may sit equally well under another concept --- that weakens
%     discrimination between items.
%   No placeholder nouns: not "the thing", "the material", "the content".
%     Name what it is: "the topic", "the activity being discussed".
%
\footnotesize
\begin{itemize}
  \item Teacher \ldots{} \emph{(an observable move, specific enough to point at on the recording)}
  \item Teacher \ldots{} \emph{(a different move --- not a restatement of the one above)}
  \item Teacher \ldots{} \emph{(name the object: the topic, the activity, the method --- never ``the thing'')}
  \item Teacher \ldots{} \emph{(nothing that would sit equally well under another concept)}
  \item Teacher \ldots{} \emph{(nothing implying a requirement the anchors do not impose)}
  \item Teacher \ldots{} \emph{(six to eight items in total)}
\end{itemize}
\normalsize
}
 
%
% SPECIAL NOTE  --- under 45 words. One or two sentences, as in TIPPS.
%   CARRIES: the exclusions most likely to produce a false positive;
%     two-directional caveats ("this is not evidence, but it is not a fault");
%     scoring allowances ("distribution across the lesson is not required").
%   DOES NOT CARRY: long exclusion lists --- most exclusions describe Column A
%     and belong in that anchor or its examples, where the coder meets them
%     while deciding; behaviours, which belong in the indicators; edge-case
%     rules, which belong in the anchor as a qualifying clause; provenance or
%     design history, which belong in the analysis document and never here.
%
\specialnote{One or two sentences, under forty-five words, holding only the exclusions most likely to be
mistaken for the practice and any caveat that cuts both ways.}
 
%
% SCORING TABLE
%   STEMS: the concept statement negated, then affirmed. Nothing else.
%
%   ANCHORS: reuse the source rubric's own wording wherever it is already
%     clear --- paraphrase is where construct drift enters.
%     NEVER refer to another column or option. Not "separates this from
%       Very Accurate", not "with limited explanation it is Column A".
%       If the substance matters, move it to the Special Note as an exclusion.
%     Read the four anchors as a sequence before shipping: adjacent options
%       must not describe the same reach or the same frequency.
%     Keep the source's vague quantifiers --- "several", "a few", "isolated",
%       "minimal" --- and do not convert them into numbers.
%     For frequency, use the TIPPS ladder: "minimal instances" ->
%       "occasionally" / "intermittent" -> "regularly"; and for a
%       more-than-once threshold, "isolated instances" versus "on more than
%       one occasion".
%     For reach, use the shared four-band scale from the front matter.
%     Prefer the plain word. If a term would make a first-week coder pause,
%       it costs more than it earns.
%     No gendered pronouns.
%
\scoregrid{Teacher does not \ldots{} \emph{(the concept statement, negated).}}%
          {Teacher \ldots{} \emph{(the concept statement, affirmed).}}{%
\footnotesize \emph{Column A, Very Accurate.} The practice is absent. Describe what its complete absence
looks like, including the form that is most often mistaken for the practice. State any condition that puts a
lesson here regardless of what else occurred. &
\footnotesize \emph{Column A, Somewhat Accurate.} The practice is attempted but does not succeed, or occurs
in a form too weak or too narrow to count. Say what is missing. Do not name another column. &
\footnotesize \emph{Column B, Somewhat Accurate.} The practice occurs genuinely but is limited --- in reach,
in frequency, or in depth. Say which. Use the shared reach bands or the TIPPS frequency ladder, not numbers. &
\footnotesize \emph{Column B, Very Accurate.} The fullest form of the practice. Describe it positively and
completely, without reference to the option beside it, and name the strongest form the source identifies.
}{%
%
% EXAMPLES  --- two or three per column. Three is a ceiling, not a quota.
%   Every example must be traceable to the source rubric. Never invent one to
%     fill a slot; two good examples beat three where the third is vague.
%   Each example in a column should show a DIFFERENT route into that option.
%   Write so that someone who has never read the rubric can follow it: name
%     the behaviour rather than gesturing at it, and name the actor rather
%     than leaving the subject implied.
%   Describe what teachers on these recordings actually do --- for instance,
%     teachers step out and return; they do not leave permanently.
%   IF THE SOURCE RUNS ONE SCENARIO ACROSS ALL FOUR OPTIONS, use it as the
%     FIRST example in every column. Holding the content constant while the
%     teaching varies is the most useful training device on the page.
%     (Concept 4: the umbrella. 5: double-digit multiplication. 6: angles.
%      7: the quiz.)
%
\exlistsize\begin{exlist}
  \item \emph{First example --- the shared scenario at this level, if the source runs one across all four.}
  \item \emph{Second example --- a different route into this option.}
  \item \emph{Third example, only if it is source-backed and adds something.}
\end{exlist} &
\exlistsize\begin{exlist}
  \item \emph{The shared scenario, one step further on.}
  \item \emph{A different route into this option.}
  \item \emph{Third only if source-backed.}
\end{exlist} &
\exlistsize\begin{exlist}
  \item \emph{The shared scenario, one step further on.}
  \item \emph{A different route into this option.}
  \item \emph{Third only if source-backed.}
\end{exlist} &
\exlistsize\begin{exlist}
  \item \emph{The shared scenario at its fullest.}
  \item \emph{A different route into this option.}
  \item \emph{Third only if source-backed.}
\end{exlist}
}
% ========================================================= END TEMPLATE ===
 
\clearpage
 
% ---------------------------------------------------------------------
%  AUTHORING CHECKLIST --- delete this page from the assembled manual.
% ---------------------------------------------------------------------
\conceptheader{}{Authoring checklist}
 
\footnotesize
{\bfseries Against the source}\\[1mm]
Read the whole guidance cell end to end before writing, not only the anchors --- a one-sentence definition
that excludes what the source includes is a construct error, not a compression. Confirm every condition in
the source is either present or deliberately cut for a stated reason. Reuse the source's wording for anchors;
paraphrase is where drift enters. Keep its vague quantifiers as they are. Check whether the original TIPPS
carries this concept and, if so, start from its paragraph. Add no rationale the source does not contain.
 
\vspace{2.5mm}
{\bfseries Slot discipline}\\[1mm]
Construct and its decisive test in the paragraph; observable behaviours in the indicators; exclusions and
two-directional caveats in the Special Note, under forty-five words; gradations in the anchors; concrete
illustrations in the examples. Nothing in the wrong slot, and nothing duplicated between the paragraph and
the note --- that duplication is usually the cheapest line to recover when a page overruns.
 
\vspace{2.5mm}
{\bfseries Language}\\[1mm]
No gendered pronouns anywhere on the page. No scoring-apparatus words in the paragraph. No references from
one anchor to another. No informal connectives such as ``Also here''. No placeholder nouns. No compressed
metaphor and no coined shorthand left undefined; if a phrase is worth keeping, spell out the mechanism beside
it. Prefer the plain word throughout.
 
\vspace{2.5mm}
{\bfseries Read it as a coder would}\\[1mm]
Read the four anchors straight through as a sequence: adjacent options must not describe the same reach or
frequency, and each must stand alone without the one beside it. Read each example as though you had never
seen the rubric. Ask of every indicator whether it would sit equally well under a different concept.
 
\vspace{2.5mm}
{\bfseries Layout}\\[1mm]
The preamble must stay byte-identical across all concepts so the pages merge cleanly. Fix overruns by
trimming text, never by changing the geometry. Working budget: paragraph up to about 190 words, six to eight
indicators, Special Note under forty-five words, source-length anchors, and three examples per column ---
four in one column will not fit alongside a long anchor row. Verify with \texttt{pdftotext -bbox} that the top
matter plus the table stays within \texttt{\textbackslash textheight} (771\,pt), and confirm the concept
occupies exactly one page.
\normalsize

 
\end{document}